
\medskip
\phantomsection\label{ctr:sga3-VIB-5.6.4}%
\noindent\textbf{Counter\SGARefLink{ctr:sga3-VIB-5.6.4}{example 5.6.4}.---}
\emph{Not every \(S\)-group \(G\) is separated. Let \(S\) be a
scheme having a non-isolated closed point \(s\), and let \(G\) be
the scheme obtained by gluing two copies of \(S\) along the open
subset \(S-\{s\}\). It is easily seen that \(G\) is not separated
over \(S\), and that it carries a natural \(S\)-group structure
whose fibers are all trivial, except for the fiber \(G_s\), which
is isomorphic to the two-element group
\(\mathbb Z/2\mathbb Z\).}%
\footnote{\label{note:sga3-VIB-45}Editor's note: we reproduced this
example in VI\(_A\), \SGARefLink{sga3:VIA:0.3}{0.3}, editor's \SGARefLink{sec:sga3-VIB-5}{note 5}.}

\medskip
\phantomsection\label{thm:sga3-VIB-5.7}%
\noindent\textbf{Theorem 5.7.---}%
\footnote{\label{note:sga3-VIB-46}Editor's note: on the one hand,
we omitted Corollary 5.6.5, which repeated \SGARefLink{cor:sga3-VIB-5.5}{5.5}. On the other hand,
the original stated \hyperref[cor:sga3-VIB-5.7.1]{Corollary 5.7.1} below as \hyperref[cor:sga3-VIB-5.7.1]{Remark 5.7}, referring
for its proof to \SGARefLink{par:sga3-VIA-4.7}{4.7}, a number that does not exist in
\emph{Lecture Notes in Mathematics} 151 (but did appear in the
1965 edition of SGAD, whose numbers \SGARefLink{par:sga3-VIA-4.5}{4.5} and \SGARefLink{par:sga3-VIA-4.6}{4.6} became \SGARefLink{lem:sga3-VIB-5.6.1}{5.6.1} and
\SGARefLink{cor:sga3-VIB-5.6.2}{5.6.2}); we have added \hyperref[thm:sga3-VIB-5.7]{Theorem 5.7}, communicated by O.\ Gabber,
which makes the aforementioned statement of SGAD more precise.}
\emph{Let \(S\) be a scheme, let \(G\) be an \(S\)-group locally
of finite presentation over \(S\), such that the function
\[
  s\longmapsto\dim G_s
\]
is locally constant on \(S\), let \(X\) be an \(S\)-scheme on
which \(G\) acts, and let \(U\) be an open subset of \(X\),
separated over \(S\). Then
\(\underline{G^0}\mathbin{\cdot}U\) is an open subset of \(X\),
separated over \(S\).}

\emph{Proof.}
Let
\[
  \mu:G\times_SX\longrightarrow X
\]
denote the action of \(G\) on \(X\). It is the composite of the
automorphism
\[
  (g,x)\longmapsto(g,gx)
\]
of \(G\times_SX\) and the projection onto \(X\). Since
\(G\times_SU\) is an open subset of \(G\times_SX\), it follows
from \hyperref[concl:sga3-VIB-4.2-iii]{4.2}\textup{(iii)} that
\[
  V=\underline{G^0}\mathbin{\cdot}U
\]
is open in \(X\).

Then, arguing as in the proof of \hyperref[cor:sga3-VIB-5.6.2]{5.6.2}, we deduce from the
implication \textup{(iv)}\(\Rightarrow\)\textup{(i)} of \hyperref[lem:sga3-VIB-5.6.3]{5.6.3}
that \(V\) is separated over \(S\).
